% -------------------------------------------------------------------
% Replacement for the \authortodo in Section III-D1
% ("Binding One Suite to Different Implementations").
%
% Describes the binding mechanism that exists in the replication
% package (oracle/conftest.py, analysis/correctness.py,
% oracle/django/conftest.py). It deliberately does NOT claim a
% specification-level name adapter: no such adapter is implemented.
% See docs/manuscript-audit.md, section 5, before the agentic
% subjects are scored for RQ1.
% -------------------------------------------------------------------

\subsubsection{Binding One Suite to Different Implementations}

Because the specification fixes neither package names nor module
layout, each subject exposes the specified behavior under a different
import path. One suite is therefore bound to each subject at run time
rather than being rewritten per subject. Every oracle test reaches the
system under test through a single canonical handle, \texttt{sut}. The
correctness runner resolves, for each subject, the pair (top-level
package name, \texttt{sys.path} entry) from the source root recorded in
the subject registry: if the root is itself a package, the import name
is the root's own name and the path entry is its \emph{parent}, never
the package directory, which would otherwise shadow standard library
modules of the same name; if the root is a container directory, the
import name is its single contained package and the path entry is the
root. The runner then prepends the path entry to \texttt{PYTHONPATH},
exports the import name in the environment, and invokes
\texttt{pytest} once per subject in a separate process; a shared
\texttt{conftest.py} imports that package and aliases it to
\texttt{sut}. An import failure is reported as a skip with the
underlying exception rather than as a behavioral failure, so a subject
that cannot be imported is never silently scored as incorrect.

Results are read back from the JUnit XML report and bucketed into the
behavioral categories of Section~\ref{sec:oracle} by test-file stem, so
the per-category tables are produced by the same mechanism for every
subject.

The Django URL module needs a configured settings module before most of
its symbols can be imported, so the Django oracle configures a minimal
in-memory settings object once per test session, with templating
enabled and no installed applications, database, or middleware, and
then calls the subject's \texttt{setup()} entry point. Requests are
constructed through the framework's own request-factory surface. This
hosts the routing and view-dispatch subsystem in isolation, without a
project skeleton on disk, so the analyzed scope and the tested scope
are identical.
